\section{Conclusion}
\label{sec:conclusion}


In this paper, we introduced \sysname{}, a framework that enables Large Language Models to analyze network configurations with the nuanced, context-aware reasoning of a human expert. We demonstrated that naive LLM applications fail due to context overload and fragmentation, and we solve this by introducing a two-part mechanism. First, a systematic context mining engine extracts all potentially relevant neighboring, similar, and referenced (both intra--- and cross---device) configurations. Second, an interactive prompting engine engages the LLM in a dialogue, allowing the model to request the specific context it needs for a final, focused analysis. Our evaluation shows that \sysname{} significantly outperforms existing tools, achieves near-perfect accuracy on known errors, and uncovers dozens of previously unknown, real-world misconfigurations. More importantly, our work reveals a critical precision-recall trade-off between different LLM backends, providing a deeper understanding of how these models can be effectively and reliably deployed for network verification.

% In this paper, we introduced \sysname{}, a framework for routing device misconfiguration detection that leverages context-aware iterative prompting to improve LLM accuracy. By systematically extracting relevant context from network configurations
% and following a guided, interactive prompting mechanism, \sysname{} addresses limitations in current model checkers, consistency checkers, and LLM-based models. Through case studies, we demonstrated that \sysname{} outperforms existing methods across different misconfiguration types, ensuring more reliable and scalable network configuration verification.